\chapter{Voorwoord}

Dit boek heeft een eenvoudige ambitie: één samenhangend biologieboek dat
alles bestrijkt van de kleuterschool tot het einde van de bachelor,
geschreven in het Nederlands, voor lezers waar ook ter wereld.

De stijl is beknopt en streng: lesteksten opgebouwd uit definities,
voorbeelden, proposities, methoden en schema's, met zorgvuldige
redeneringen zodra die op het betreffende niveau toegankelijk zijn,
gevolgd door zorgvuldig gekozen oefeningen. Resultaten die zonder
volledige afleiding worden gegeven, zijn uitdrukkelijk als toegegeven
gemarkeerd. Van alle oefeningen staan achter in het boek volledige
uitwerkingen. De biologiehoofdstukken geven de voorkeur aan figuren boven
formules.

Elk schooljaar vormt een deel van het boek, en elk deel is in hoofdstukken
verdeeld. Hoe vroeger het jaar, hoe jonger de lezer voor wie het geschreven
is: meer illustraties, meer uitgewerkte tussenstappen en mildere
oefeningen.

\section*{Hoe je dit boek leest}

De uitspraken worden per hoofdstuk genummerd en zijn met kleur gecodeerd:
definities in blauw, stellingen in rood, proposities en verwanten in
oranje, en methoden --- korte recepten voor standaardtaken --- in groen.
De oefeningen lopen van ($\star$) voor rechtstreekse toepassingen tot
($\star\star\star$) voor lastigere opgaven.

\section*{Meewerken}

De bron van dit boek staat in een openbare git-repository:
\begin{center}
\url{https://github.com/bvirrion/one-biology-book}
\end{center}
Bijdragen --- nieuwe hoofdstukken, correcties, betere bewijzen, extra
oefeningen --- zijn welkom; zie \texttt{CONTRIBUTING.md} in de hoofdmap van
de repository.
